\clearpage
\section{Results}
\label{sec:results}

This laboratory protocol examines the photocurrent spectrum of five different photodiodes
(D1, \dots , D5), see %TODO the dots are not formatted correctly
\cref{subsec:current_spectrum}. By recording a reference spectrum in \cref{subsec:reference_spectra},
the spectral responsivity and quantum efficiency of those photodiodes can be
extracted, see \cref{subsec:spectral_responsivity}. The response linearity with
respect to the optical power is another key metric that is examined in
\cref{subsec:linearity}. Although D3, an ordinary silicon photodiode, and D1, a
uv-enhanced silicon photodiode, are based on silicon, their responsivity
spectrum differ significantly. This difference is discussed in \cref{subsec:comparison}.
D4's responsivity is compared to a human eye response curve $V(\lambda )$, see \cref{subsec:eye}.
The response of the diode depends on the frequency of the incoming radiation. This
was examined in \cref{subsec:chopper} by varying the reference frequency and therefore
changing the angular speed of the chopper wheel.

\subsection{Diode Current Spectrum}
\label{subsec:current_spectrum}

\begin{figure}
	\centering
	\includegraphics{../plots/photocurrent.pdf}
	\caption{Photocurrent $I_{\mathrm{ph}}$ of the five photodiodes as a function
	of wavelength $\lambda$. Markers indicate the maximum of each spectrum at $\lambda
	_{\mathrm{max}}$.}
	\label{fig:photocurrent}
\end{figure}

Photocurrent spectra of five photodiodes were recorded and are depicted in \cref{fig:photocurrent}.
As all photodiodes are silicon-based, the band gap of \qty{1.124}{\electronvolt} determines
a minimal energy to excite electrons. This corresponds to a threshhold
wavelength of
$\lambda_{\mathrm{th}}=hc / E_{\mathrm{g}}=\qty{1103}{\nano\meter}$. Above this threshhold
wavelength, no signal intensity is measured.

The photocurrent-spectra can be grouped into two categories, based on their line profile:
D3 and D5 peak at \qty{1000}{\nano \meter} while D1, D2 and D4 at around \qty{600}{\nano\meter}.
This difference is discussed in \cref{subsec:comparison}. Te photocurrent significantly
drops for small wavelengths. This is a direct consequence of the decrease in
optical power of the halogen lamp, see \cref{subsec:reference_spectra}.

\subsection{Reference Spectra}
\label{subsec:reference_spectra}

\begin{figure}
	\centering
	\begin{subfigure}
		{0.48\linewidth}
		\centering
		\includegraphics[width=\linewidth]{../plots/reference_spectra.pdf}
		\caption{Linear scale.}
		\label{fig:reference_spectra}
	\end{subfigure}
	\hfill
	\begin{subfigure}
		{0.48\linewidth}
		\centering
		\includegraphics[width=\linewidth]{../plots/reference_spectra_log.pdf}
		\caption{Logarithmic scale.}
		\label{fig:reference_spectra_log}
	\end{subfigure}
	\caption{Optical power $P_{\mathrm{opt}}$ of the reference spectra recorded at
	the beginning and at the end of the experiment, together with their average.}
	\label{fig:reference}
\end{figure}

\begin{figure}
	\centering
	\includegraphics{../plots/reference_spread.pdf}
	\caption{Relative deviation between the two reference spectra with respect to
	their average.}
	\label{fig:reference_spread}
\end{figure}

A calibrated reference detector with single-point calibration at \qty{633}{\nano\meter}
for radiometric power measurement was used to determine the optical reference
power spectrum. To capture the thermal shifts of the halogen-lamp, the spectrum
was recorded at the beginning and at the end of the experiment, see \cref{fig:reference}.
As the lamp was operating for several hours prior to the beginning of the
experiment, only a small shift is observed. Time-dependent effects can be
neglected.

The relative deviation of the optical power is shown in
\cref{fig:reference_spread}. As the intensity drops below \qty{550}{\nano\meter},
the signal enters the noise range. Subsequent procedures to identify the maximum of
the spectrum only consider wavelengths greater tha \qty{550}{\nano\meter}.

\subsection{Spectral Responsivity and Quantum Efficiency}
\label{subsec:spectral_responsivity}

%TODO These two graphs should be plotted next to each other using minipage i guess?. They should render fine without rescaling explicitely (but this may not be the case)
\begin{figure}
	\centering
	\includegraphics{../plots/responsivity.pdf}
	\caption{Spectral responsivity $R_{\lambda}$ of the five photodiodes.}
	\label{fig:responsivity}
\end{figure}

\begin{figure}
	\centering
	\includegraphics{../plots/quantum_efficiency.pdf}
	\caption{External quantum efficiency $\eta$ of the five photodiodes as a
	function of wavelength $\lambda$.}
	\label{fig:quantum_efficiency}
\end{figure}

Using the average reference spectrum $P_{\mathrm{opt}}^{\mathrm{avg}}$, the responsivity
$R_{\lambda}=I_{\mathrm{ph}}/ P_{\mathrm{opt}}^{\mathrm{avg}}$ was calculated and
is depicted in \cref{fig:responsivity}. $R_{\lambda}$ does not depend on the
optical sources power and is therefore a better device characteristic. The band gap
still determines the highest wavelength threshold. The maximum responsivity
wavelengths for each diodes are depicted in \cref{fig:responsivity}, they are similar
to raw maxima. %TODO This sounds really bad
The steep decrease towards low wavelengths is corrected, as it was caused by the
vanishing detector power.

The external quantum efficiency spectrum is plotted in
\cref{fig:quantum_efficiency}. D3 features the highest quantum efficiency maximum
at \qty{37.14}{\percent} while D5 features the lowest at \qty{2.32}{\percent}.
The characteristic values of all photodiodes are summarized in \cref{tab:diodes}.

Direct comparisons between the responsivities and quantum efficiencies are not
meaningful, given the high systematic error in the measurement setup. The
photodiode mount was readjusted after each photodiode installation. The detection
area generally differs from diode to diode as well as in comparison to the detector.
Therefore, the responsivity and external quantum efficiency is only a rough estimate. %TODO nur die relativen Werte (normiert auf den Maximalwert) lassen sich gut vergleichen

\begin{table}
	\centering
	\begin{tabular}{lSSSSS}
\toprule
{} & {\CellWithForceBreak{$\lambda_\mathrm{max}$\\{}[\unit{\nano\meter}]}} & {\CellWithForceBreak{$R(\lambda_\mathrm{max})$\\{}[\unit{\ampere\per\watt}]}} & {\CellWithForceBreak{$\eta(\lambda_\mathrm{max})$\\{}[\unit{\percent}]}} & {\CellWithForceBreak{$R(\qty{550}{\nano\meter})$\\{}[\unit{\ampere\per\watt}]}} & {\CellWithForceBreak{$\eta(\qty{550}{\nano\meter})$\\{}[\unit{\percent}]}} \\
\midrule
Diode 1 & 558 & 0.0774 & 17.17 & 0.0767 & 17.25 \\
Diode 2 & 558 & 0.0570 & 12.66 & 0.0567 & 12.75 \\
Diode 3 & 898 & 0.2692 & 37.14 & 0.1585 & 35.67 \\
Diode 4 & 582 & 0.0128 & 2.73 & 0.0121 & 2.73 \\
Diode 5 & 939 & 0.0176 & 2.32 & 0.0018 & 0.41 \\
\bottomrule
\end{tabular}

	\caption{Maximum responsivity $R$ and external quantum efficiency $\eta$ of
	each photodiode together with the wavelength $\lambda_{\mathrm{max}}$ at which
	they occur, compared to the values at the fixed wavelength \qty{550}{\nano\meter}.}
	\label{tab:diodes}
\end{table}

\subsection{Photodiode Detector Linearity}
\label{subsec:linearity}

\begin{figure}
	\centering
	\includegraphics{../plots/linearity.pdf}
	\caption{Relative photocurrent of diodes 4 and 5 as a function of the relative
	incident photon flux, together with power-law fits of slope $m$.}
	\label{fig:linearity}
\end{figure}

The photocurrents of D4 and D5 were measured at a fixed wavelength while the optical
power was variied using neutral density filters, depicted in
\cref{fig:linearity}. As the theory suggests $I_{\mathrm{ph}}\propto P_{\mathrm{opt}}$,
the following relations holds true:
\begin{equation}
	\frac{I}{I_{0}}=\frac{P}{P_{0}}.
\end{equation}
The slope $m$ of a linear function $f(x)=mx+n$ fitted on the logarithmic data indicates
the exponent as $f(x)=\exp(n)x^{m}$ of the original data. %TODO: Please refactor this sentence,

D4 shows a nearly linear scaling as $m=0.92$. However, D5 shows two regimes, with
an unphysical break at \num{2e-2}. While the second regime features a nearly linear
scaling with $m=0.94$, the first regime only scales as $m=0.62$.

\subsection{Comparison between Si- and UV-Enhanced Si- Photodiodes}
\label{subsec:comparison}

Both, the Si-Photodiode D3 as well as the uv-enhanced Si-Photodiode D1 are
silicon based and therefore share the threshhold wavelength of \qty{1103}{\nano\meter}.
However, their line profiles differ significantly. D1 has a maximum responsivity at
\qty{558}{\nano\meter}, while D3 features a maxmimum responsivity at \qty{898}{\nano\meter}.

As the measurement does not enter the UV-regime (\qtyrange{250}{400}{\nano\meter}),
no statements can be concluded about their responsivity in this region. It is expected
that D1 features a significantly higher responsivity.

For an ordinary silicon photodiode, infrared radiation penetrates deeply into
the material due to silicon's reduced absorption coefficient: The generation of electron-hole
pairs takes place in the whole SCR. However UV radiation does not penetrate deep
into the material due to a higher absorption coefficient. A large proportion of electron-hole
pair generation happens in the inversion layer. These pairs have a high chance
to recombine at the surface and effectively do not contribute to the
photocurrent. The Schottky photodiode therefore features a responsivity falloff in
the UV-regime attributed to the presence of an inversion layer in the depletion region.

UV-enhanced Si-photodiodes may feature higher doping concentrations $N_{\mathrm{D}}$.
This implies a shrinking inversion layer width as well as a shrinking space
charge region width as both scale as $\propto 1 / \sqrt{N_{\mathrm{D}}}$
\autocite{korde}. UV radiation is mostly absorbed in the SCR and electron-hole
excitation results in a current. However the SCR is too narrow for infrared radiation,
as a large proportion passes through it. This optimization leads to a responsivity
loss in the near infrared regime.

Different coating may additionally influence the responsivity as well. By
comparing line profiles, D5 could also be an ordinary photodiode, while D2 and
D4 might be uv-enhanced photodiodes.

\subsection{Human Eye Response}
\label{subsec:eye}

\begin{figure}
	\centering
	\includegraphics{../plots/human_eye.pdf}
	\caption{Normalized spectral responsivity of diode 4 compared to the photopic
	and scotopic response curves $V(\lambda)$ of the human eye.}
	\label{fig:human_eye}
\end{figure}

This section is based on \citeauthoryear{paschotta}. Human eye vision is based
on four types of photorecepters: Rods as well as three types on cones. Depending
on the incoming luminance, the brain uses different photoreceptor signals to determine
the brightness. Under bright-light conditions, the brain selects photopic vision,
while under low-light conditions, scotopic vision is activated. Each
photorecepter uses a different photopigment that changes chemically under light
absorption.

Scotopic vision uses rods. The rhodopsin photopigment is more light sensitive than
cone pigments. Photopic vision uses cones. Each cone type features its own photopigment
with sligthly different absorption maxima.

The $V(\lambda)$ curves for scotopic and photopic vision, depicted together with
the normalized responsivity of D4 in \cref{fig:human_eye}, are experimentally determined.
It is not a physical property of a photoreceptor, but the outcome of a psychological
exeriment averaged over a panel of human observers \autocite{grundmann}. Each observer
equalizes the brightness of a monochromatic light with $\lambda_{0}= \qty{555}{\nano\meter}$
to that of a light source with wavelength $\lambda$ with the same radiation power
using flicker photometry. The experiment was conducted under bright-light and
under low-light conditions and the relative radiative power is used to construct
$V(\lambda)$.

The three curves feature a similar maximum, at around $\lambda = \qty{550}{\nano\meter}$.
However, the peak width of D4 is significantly larger than both $V(\lambda)$ curves.
Without any correction filter, D4 cannot be used as a photometric diode.

\subsection{Chopper Frequency Dependence}
\label{subsec:chopper}

\begin{figure}
	\centering
	\includegraphics{../plots/chopper.pdf}
	\caption{Photocurrent $I_{\mathrm{ph}}$ as a function of the reference
	frequency $f$. Blue/green/orange dots indicate an integration time of \qty{1}{\second}/
	\qty{2}{\second} /\qty{10}{\second}}
	. \label{fig:chopper}
\end{figure}

A frequency dependence for a diode on the intensity is expected. In
order to observe this effect, the photocurrent was measured at a fixed wavelength
for different reference frequencies, resulting in differen chopper angular
velocitie. The result is depicted in \cref{fig:chopper}. No change in
photocurrent is observed for high frequencies.

However, a sharp falloff at low frequencies is present. A possible explanation
of this behaviour is a coupling condensator that connects two stages acting as a
high-pass filter. Using the fit function
\begin{equation}
	I_{\mathrm{ph}}(f)=\frac{I_{\mathrm{\mathrm{sat}}}f}{\sqrt{ f^{2}+f_{\mathrm{c}}^{2}}}
\end{equation}

with the saturation current $I_{\mathrm{sat}}$ and the cutoff frequency
$f_{\mathrm{c}}$, the fit yields $f_{\mathrm{c}}=\qty{0.77}{\hertz}$ and $I_{\mathrm{sat}}
=\qty{125}{\nano\ampere}$.